%% ======================================================================
%% DROPPED 2026-08-27, OWNER'S EXPLICIT INSTRUCTION: "section 8.6, two pages or
%% fewer".  THIS FIGURE IS NO LONGER INPUT BY ANYTHING.  Section 8.6 now has no
%% figure: two prose fragments and the inline table tab_eis_codelevel.
%% NOTHING BELOW THIS HEADER WAS EDITED.  Restoring the figure is two
%% mechanical steps in note_eis_intro.tex: put
%%   %% ======================================================================
%% DROPPED 2026-08-27, OWNER'S EXPLICIT INSTRUCTION: "section 8.6, two pages or
%% fewer".  THIS FIGURE IS NO LONGER INPUT BY ANYTHING.  Section 8.6 now has no
%% figure: two prose fragments and the inline table tab_eis_codelevel.
%% NOTHING BELOW THIS HEADER WAS EDITED.  Restoring the figure is two
%% mechanical steps in note_eis_intro.tex: put
%%   %% ======================================================================
%% DROPPED 2026-08-27, OWNER'S EXPLICIT INSTRUCTION: "section 8.6, two pages or
%% fewer".  THIS FIGURE IS NO LONGER INPUT BY ANYTHING.  Section 8.6 now has no
%% figure: two prose fragments and the inline table tab_eis_codelevel.
%% NOTHING BELOW THIS HEADER WAS EDITED.  Restoring the figure is two
%% mechanical steps in note_eis_intro.tex: put
%%   %% ======================================================================
%% DROPPED 2026-08-27, OWNER'S EXPLICIT INSTRUCTION: "section 8.6, two pages or
%% fewer".  THIS FIGURE IS NO LONGER INPUT BY ANYTHING.  Section 8.6 now has no
%% figure: two prose fragments and the inline table tab_eis_codelevel.
%% NOTHING BELOW THIS HEADER WAS EDITED.  Restoring the figure is two
%% mechanical steps in note_eis_intro.tex: put
%%   \input{include/analog/fig_eis_nyquist.tex}
%% back after the conditions paragraph, and put "(Figure~\ref{fig:eis-nyquist})"
%% back on the sentence that now ends "Sweeps carry 20 points per decade."
%% ARCHIVE that holds this figure and every datum in it:
%%   ~/chips/castalia/simruns/eisfigs2_20260826  (build/fig1_nyquist.tex,
%%   data/fig1_{fresh,typical,fouled}_{true,meas,mark}.dat), reduced from
%%   ~/chips/castalia/simruns/eisdeep_20260826 csv/nyquist_*_n0.csv.
%% The caption's facts and the disposition of each are written out fact by fact
%% in note_eis_intro.tex's header, and summarised in note_eis_accuracy.tex's.
%% In short: the agreement figures and the gain-code-0 coverage limit duplicate
%% or are arithmetic on headlines the prose still carries; the panel scales,
%% the equal axes and the "fresh arm above 151 kOhm" note are figure-reading
%% caveats with no figure to read; and the one number that lived only here --
%% 2.94 % and 4.67 deg at 100 kHz -- was CUT rather than moved, because it is
%% the far end of this figure's own code-0 envelope and a decade above the
%% usable band the section states.  It is in the archive above.
%% ======================================================================
%%
%% HAND-WRITTEN fragment -- NOT generated by maestro2tex (raw Spectre runs, not
%% a Maestro results database; there is no maestro2tex config for any *_eis_*
%% fragment).  Precedent for the style block and \MaestroRoot convention:
%% fig_pixtop_noise.tex, fig_eis_error.tex.
%%
%% Drop-in source: ~/chips/castalia/simruns/eisfigs2_20260826/build/
%% fig1_nyquist.tex, with the standalone preamble, the in-figure legend row and
%% the in-figure note block removed -- the note is the caption here.
%% Data: data/eis_nyq_{fresh,typical,fouled}_{true,meas,dec}.dat, copied from
%% eisfigs2_20260826/data/fig1_<load>_{true,meas,mark}.dat, reduced there from
%% ~/chips/castalia/simruns/eisdeep_20260826/csv/nyquist_{fresh,typical,
%% fouled}_n0.csv -- three of that campaign's 30 ac_dense cells, 143
%% frequencies each, 20 points/decade from 0.01 Hz, columns ReZ_meas/ImZ_meas
%% (architecture A) and ReZ_true/ImZ_true (the closed-form Randles impedance)
%% taken as written.  Run provenance: ac_dense/ac_results.json, csv/index.json.
%% Units are ohms in the .dat; the axes divide by 1e6 (fresh) or 1e3.
%%
%% EQUAL AXES ARE LOAD-BEARING.  Every panel has height = 0.55 x width and
%% limits set to the same 0.55 ratio, so the arc is geometrically a semicircle
%% and a departure from a circle is a datum.  Do not change one of width,
%% height or the limits without changing all three.  All 143 measured points
%% are plotted -- the density IS the result, so no mark thinning here (unlike
%% fig_eis_error, where the curve and not the grid is the subject).
%%
%% Key numbers (eisfigs2 README S1): chord deviation of the plotted polygon
%% from the true arc, as a fraction of the arc radius and identical for all
%% three loads because the shape is -- 0.165 % at 20 points/decade, 0.651 % at
%% the archived 10, 0.074 % at 30, against a 0.30 % instrument error.  |D|Z||
%% fresh 0.294 / typical 0.423 / fouled 0.436 % worst case below 20 kHz and
%% 2.808 / 2.925 / 2.936 % at 100 kHz; phase <= 1.031 deg below 20 kHz and
%% <= 4.667 deg at 100 kHz.  All three arcs peak at 1.59 Hz (R_ct C_dl = 0.1 s).
%% Gain code N = 0, R_f = 19.63 kOhm measured in the closed loop.
%% Caveats that had to survive into the caption: the panels have INDEPENDENT
%% scales and equal axes; the ac analysis carries no converter, so the fresh
%% arm above |Z|max = 150.9 kOhm at N = 0 would be taken at N = 20 in firmware
%% (coverage_map.txt).  Nominal corner castalia_typ_25, 37 C, no mismatch -- the
%% corner box on this quantity is 1.6-1.8x (eisdeep S5.3) and the Monte Carlo
%% spread on a stored-law reading is 10 % (~/chips/castalia/simruns/
%% eisdeep_20260826 mc_report/; fig_eis_mccal.tex plotted it until the second
%% cut of 2026-08-26 de-inputted it, and it is still on disk).
%% Not drawn: corner overlays.  The 120 dense corner spectra exist (csv/corner_*)
%% but 5 corners x 3 panels destroys the one thing this figure is for.
%% Library cell names are kept out of the rendered prose.
%% 2026-08-26: the Table~\ref{tab:eis-coverage} pointer was dropped from the
%% caption -- that table went with the cut of note_eis_coverage.tex.  The 236 ohm
%% - 150.9 kOhm window it named is stated inline here and its source is
%% ~/chips/castalia/simruns/eisdemo_20260825/coverage_map.txt.
%% 2026-08-27: the caption was TIGHTENED (wording only -- no number, condition
%% or caveat removed).  Cut: "of them" after all 143; "the loads span" before
%% the three decade scales; "a semicircle reads as a semicircle and" (the
%% second half of the clause says it); "no longer"/"the archived" as filler.
%% "the current measured" became "the current at the transimpedance output",
%% which is the architecture the run used.  Done with the same pass that
%% replaced fig_eis_codelevel with tab_eis_codelevel.
%%
%% 2026-08-27, SECOND PASS -- CONTENT CUT, NOT WORDING.  The earlier pass above
%% was wording only and left the caption at 5 sentences but TWELVE rendered
%% lines (163 pt in the built PDF), which is over the chapter's 2-5 sentence
%% caption style by volume even though it met it by count.  Now 4 sentences.
%% WHAT LEFT THE CAPTION, and where it still lives:
%%   - THE DENSITY CLAIM AND THE CHORD-DEVIATION RESULT, cut together as a
%%     unit: "all 143" of the plotted points, the 0.165 % chord deviation of
%%     the plotted polygon from the true arc, the 0.651 % at ten points per
%%     decade and the 0.30 % instrument error against which both were read.
%%     The rule is that a number keeps its qualifier or goes with it: the
%%     0.165 % only means anything beside the instrument error, and the
%%     instrument error only means anything beside the density claim it
%%     licenses, so all four go together rather than any one standing bare.
%%     ARCHIVE: ~/chips/castalia/simruns/eisfigs2_20260826, README S1 (and the
%%     0.074 % at 30 points/decade, which was never in the caption).  All 143
%%     points ARE still plotted -- no mark thinning -- see the note above.
%%   - "with the applied potential read at the reference electrode and the
%%     current at the transimpedance output": the sensing architecture, still
%%     stated in the section's own opening paragraph ("A1 applies it at the
%%     reference electrode, A2 and the converter read the resulting current
%%     through the same programmable feedback resistor"), so it is present in
%%     the section, once, and not repeated here.
%%   - the explicit \SIlist of ring frequencies -> "rings mark the decades".
%%     NOT a content cut: every ring carries its own printed label inside the
%%     panels (0.01 Hz / 0.1 Hz / 1 Hz / 10 Hz / 100 Hz nodes), so the list was
%%     the caption reading the figure aloud.
%% WHAT HAD TO STAY, and did: the INDEPENDENT panel scales and the equal axes
%% (without both, a reader misreads the arcs); the agreement figures and where
%% they degrade (0.44 % / 1.03 deg below 20 kHz -> 2.94 % / 4.67 deg at
%% 100 kHz); the gain-code coverage limit (236 ohm - 150.9 kOhm at code 0, so
%% the fresh arm above 151 kOhm is a code-20 read in firmware).
\providecommand{\MaestroRoot}{include/analog/}
\providecolor{mtxA}{HTML}{2A78D6}
\begin{figure}[htbp]
  \centering
  {\pgfplotsset{compat=1.18}%
  \pgfplotsset{
    nyqax/.style={
      scale only axis=true,
      width=3.95cm, height=2.1725cm,
      grid=both,
      major grid style={line width=0.2pt, draw=black!14},
      minor grid style={line width=0.2pt, draw=black!7},
      axis line style={line width=0.4pt, draw=black!45},
      tick style={line width=0.4pt, draw=black!45},
      tick align=outside,
      label style={font=\scriptsize}, tick label style={font=\scriptsize},
      xmin=0, clip=false,
    }}%
  \begin{tikzpicture}
    \begin{axis}[nyqax,
      xmax=1.06106, ymin=-0.0262612, ymax=0.557322,
      xtick={0,0.5,1}, ytick={0,0.2,0.4},
      xlabel={\(Z'\)~[\si{\mega\ohm}]}, ylabel={\(-Z''\)~[\si{\mega\ohm}]},
      title={\footnotesize\textbf{fresh}},
    ]
      \addplot[black!65, line width=0.7pt, no marks]
        table[x expr=\thisrow{x}/1e6, y expr=\thisrow{y}/1e6]
        {\MaestroRoot data/eis_nyq_fresh_true.dat};
      \addplot[only marks, mark=*, mark size=0.85pt, mtxA,
               mark options={fill=mtxA, draw=none}]
        table[x expr=\thisrow{x}/1e6, y expr=\thisrow{y}/1e6]
        {\MaestroRoot data/eis_nyq_fresh_meas.dat};
      \addplot[only marks, mark=o, mark size=1.9pt, black, line width=0.45pt]
        table[x expr=\thisrow{x}/1e6, y expr=\thisrow{y}/1e6]
        {\MaestroRoot data/eis_nyq_fresh_dec.dat};
      \node[left=2pt, font=\tiny] at (axis cs:1.00096,0.00628294) {\SI{0.01}{\hertz}};
      \node[above left=1pt, font=\tiny] at (axis cs:0.997068,0.0625848) {\SI{0.1}{\hertz}};
      \node[above right=1pt, font=\tiny] at (axis cs:0.717957,0.450477) {\SI{1}{\hertz}};
      \node[below right=1pt, font=\tiny] at (axis cs:0.0257045,0.155223) {\SI{10}{\hertz}};
      \node[right=2pt, font=\tiny] at (axis cs:0.00125324,0.0159115) {\SI{100}{\hertz}};
    \end{axis}
  \end{tikzpicture}\hfill
  \begin{tikzpicture}
    \begin{axis}[nyqax,
      xmax=107.06, ymin=-2.64974, ymax=56.2333,
      xtick={0,50,100}, ytick={0,20,40},
      xlabel={\(Z'\)~[\si{\kilo\ohm}]}, ylabel={\(-Z''\)~[\si{\kilo\ohm}]},
      title={\footnotesize\textbf{typical}},
    ]
      \addplot[black!65, line width=0.7pt, no marks]
        table[x expr=\thisrow{x}/1000, y expr=\thisrow{y}/1000]
        {\MaestroRoot data/eis_nyq_typical_true.dat};
      \addplot[only marks, mark=*, mark size=0.85pt, mtxA,
               mark options={fill=mtxA, draw=none}]
        table[x expr=\thisrow{x}/1000, y expr=\thisrow{y}/1000]
        {\MaestroRoot data/eis_nyq_typical_meas.dat};
      \addplot[only marks, mark=o, mark size=1.9pt, black, line width=0.45pt]
        table[x expr=\thisrow{x}/1000, y expr=\thisrow{y}/1000]
        {\MaestroRoot data/eis_nyq_typical_dec.dat};
      \node[left=2pt, font=\tiny] at (axis cs:100.996,0.628294) {\SI{0.01}{\hertz}};
      \node[above left=1pt, font=\tiny] at (axis cs:100.607,6.25848) {\SI{0.1}{\hertz}};
      \node[above right=1pt, font=\tiny] at (axis cs:72.6957,45.0477) {\SI{1}{\hertz}};
      \node[below right=1pt, font=\tiny] at (axis cs:3.47045,15.5223) {\SI{10}{\hertz}};
      \node[right=2pt, font=\tiny] at (axis cs:1.02532,1.59115) {\SI{100}{\hertz}};
    \end{axis}
  \end{tikzpicture}\hfill
  \begin{tikzpicture}
    \begin{axis}[nyqax,
      xmax=11.66, ymin=-0.288585, ymax=6.12441,
      xtick={0,5,10}, ytick={0,2,4,6},
      xlabel={\(Z'\)~[\si{\kilo\ohm}]}, ylabel={\(-Z''\)~[\si{\kilo\ohm}]},
      title={\footnotesize\textbf{fouled}},
    ]
      \addplot[black!65, line width=0.7pt, no marks]
        table[x expr=\thisrow{x}/1000, y expr=\thisrow{y}/1000]
        {\MaestroRoot data/eis_nyq_fouled_true.dat};
      \addplot[only marks, mark=*, mark size=0.85pt, mtxA,
               mark options={fill=mtxA, draw=none}]
        table[x expr=\thisrow{x}/1000, y expr=\thisrow{y}/1000]
        {\MaestroRoot data/eis_nyq_fouled_meas.dat};
      \addplot[only marks, mark=o, mark size=1.9pt, black, line width=0.45pt]
        table[x expr=\thisrow{x}/1000, y expr=\thisrow{y}/1000]
        {\MaestroRoot data/eis_nyq_fouled_dec.dat};
      \node[left=2pt, font=\tiny] at (axis cs:10.9996,0.0628294) {\SI{0.01}{\hertz}};
      \node[above left=1pt, font=\tiny] at (axis cs:10.9607,0.625848) {\SI{0.1}{\hertz}};
      \node[above right=1pt, font=\tiny] at (axis cs:8.16957,4.50477) {\SI{1}{\hertz}};
      \node[below right=1pt, font=\tiny] at (axis cs:1.24705,1.55223) {\SI{10}{\hertz}};
      \node[right=2pt, font=\tiny] at (axis cs:1.00253,0.159115) {\SI{100}{\hertz}};
    \end{axis}
  \end{tikzpicture}}
  \caption[{Nyquist arcs of the three electrode models, measured against
  analytic.}]{Nyquist arcs of the three electrode models measured by the
  channel (dots) against the analytic impedance of the load (line),
  \SI{0.01}{\hertz} to \SI{100}{\kilo\hertz} at \num{20} points per decade,
  gain code 0, nominal process point, \SI{37}{\celsius}; rings mark the
  decades. \textbf{The three panels have independent scales} ---
  \SI{1}{\mega\ohm}, \SI{100}{\kilo\ohm}, \SI{10}{\kilo\ohm} --- and each has
  equal axes, so a departure from a circle is a measurement and not an aspect
  ratio. Agreement is within \SI{0.44}{\percent} of magnitude and
  \SI{1.03}{\degree} of phase below \SI{20}{\kilo\hertz} on all three loads,
  rising to \SI{2.94}{\percent} and \SI{4.67}{\degree} at
  \SI{100}{\kilo\hertz}. Gain code 0 covers \SI{236}{\ohm} to
  \SI{150.9}{\kilo\ohm}, a limit the small-signal sweep does not have, so the
  fresh arm above \SI{151}{\kilo\ohm} is a gain-code-20 read in firmware.}
  \label{fig:eis-nyquist}
\end{figure}

%% back after the conditions paragraph, and put "(Figure~\ref{fig:eis-nyquist})"
%% back on the sentence that now ends "Sweeps carry 20 points per decade."
%% ARCHIVE that holds this figure and every datum in it:
%%   ~/chips/castalia/simruns/eisfigs2_20260826  (build/fig1_nyquist.tex,
%%   data/fig1_{fresh,typical,fouled}_{true,meas,mark}.dat), reduced from
%%   ~/chips/castalia/simruns/eisdeep_20260826 csv/nyquist_*_n0.csv.
%% The caption's facts and the disposition of each are written out fact by fact
%% in note_eis_intro.tex's header, and summarised in note_eis_accuracy.tex's.
%% In short: the agreement figures and the gain-code-0 coverage limit duplicate
%% or are arithmetic on headlines the prose still carries; the panel scales,
%% the equal axes and the "fresh arm above 151 kOhm" note are figure-reading
%% caveats with no figure to read; and the one number that lived only here --
%% 2.94 % and 4.67 deg at 100 kHz -- was CUT rather than moved, because it is
%% the far end of this figure's own code-0 envelope and a decade above the
%% usable band the section states.  It is in the archive above.
%% ======================================================================
%%
%% HAND-WRITTEN fragment -- NOT generated by maestro2tex (raw Spectre runs, not
%% a Maestro results database; there is no maestro2tex config for any *_eis_*
%% fragment).  Precedent for the style block and \MaestroRoot convention:
%% fig_pixtop_noise.tex, fig_eis_error.tex.
%%
%% Drop-in source: ~/chips/castalia/simruns/eisfigs2_20260826/build/
%% fig1_nyquist.tex, with the standalone preamble, the in-figure legend row and
%% the in-figure note block removed -- the note is the caption here.
%% Data: data/eis_nyq_{fresh,typical,fouled}_{true,meas,dec}.dat, copied from
%% eisfigs2_20260826/data/fig1_<load>_{true,meas,mark}.dat, reduced there from
%% ~/chips/castalia/simruns/eisdeep_20260826/csv/nyquist_{fresh,typical,
%% fouled}_n0.csv -- three of that campaign's 30 ac_dense cells, 143
%% frequencies each, 20 points/decade from 0.01 Hz, columns ReZ_meas/ImZ_meas
%% (architecture A) and ReZ_true/ImZ_true (the closed-form Randles impedance)
%% taken as written.  Run provenance: ac_dense/ac_results.json, csv/index.json.
%% Units are ohms in the .dat; the axes divide by 1e6 (fresh) or 1e3.
%%
%% EQUAL AXES ARE LOAD-BEARING.  Every panel has height = 0.55 x width and
%% limits set to the same 0.55 ratio, so the arc is geometrically a semicircle
%% and a departure from a circle is a datum.  Do not change one of width,
%% height or the limits without changing all three.  All 143 measured points
%% are plotted -- the density IS the result, so no mark thinning here (unlike
%% fig_eis_error, where the curve and not the grid is the subject).
%%
%% Key numbers (eisfigs2 README S1): chord deviation of the plotted polygon
%% from the true arc, as a fraction of the arc radius and identical for all
%% three loads because the shape is -- 0.165 % at 20 points/decade, 0.651 % at
%% the archived 10, 0.074 % at 30, against a 0.30 % instrument error.  |D|Z||
%% fresh 0.294 / typical 0.423 / fouled 0.436 % worst case below 20 kHz and
%% 2.808 / 2.925 / 2.936 % at 100 kHz; phase <= 1.031 deg below 20 kHz and
%% <= 4.667 deg at 100 kHz.  All three arcs peak at 1.59 Hz (R_ct C_dl = 0.1 s).
%% Gain code N = 0, R_f = 19.63 kOhm measured in the closed loop.
%% Caveats that had to survive into the caption: the panels have INDEPENDENT
%% scales and equal axes; the ac analysis carries no converter, so the fresh
%% arm above |Z|max = 150.9 kOhm at N = 0 would be taken at N = 20 in firmware
%% (coverage_map.txt).  Nominal corner castalia_typ_25, 37 C, no mismatch -- the
%% corner box on this quantity is 1.6-1.8x (eisdeep S5.3) and the Monte Carlo
%% spread on a stored-law reading is 10 % (~/chips/castalia/simruns/
%% eisdeep_20260826 mc_report/; fig_eis_mccal.tex plotted it until the second
%% cut of 2026-08-26 de-inputted it, and it is still on disk).
%% Not drawn: corner overlays.  The 120 dense corner spectra exist (csv/corner_*)
%% but 5 corners x 3 panels destroys the one thing this figure is for.
%% Library cell names are kept out of the rendered prose.
%% 2026-08-26: the Table~\ref{tab:eis-coverage} pointer was dropped from the
%% caption -- that table went with the cut of note_eis_coverage.tex.  The 236 ohm
%% - 150.9 kOhm window it named is stated inline here and its source is
%% ~/chips/castalia/simruns/eisdemo_20260825/coverage_map.txt.
%% 2026-08-27: the caption was TIGHTENED (wording only -- no number, condition
%% or caveat removed).  Cut: "of them" after all 143; "the loads span" before
%% the three decade scales; "a semicircle reads as a semicircle and" (the
%% second half of the clause says it); "no longer"/"the archived" as filler.
%% "the current measured" became "the current at the transimpedance output",
%% which is the architecture the run used.  Done with the same pass that
%% replaced fig_eis_codelevel with tab_eis_codelevel.
%%
%% 2026-08-27, SECOND PASS -- CONTENT CUT, NOT WORDING.  The earlier pass above
%% was wording only and left the caption at 5 sentences but TWELVE rendered
%% lines (163 pt in the built PDF), which is over the chapter's 2-5 sentence
%% caption style by volume even though it met it by count.  Now 4 sentences.
%% WHAT LEFT THE CAPTION, and where it still lives:
%%   - THE DENSITY CLAIM AND THE CHORD-DEVIATION RESULT, cut together as a
%%     unit: "all 143" of the plotted points, the 0.165 % chord deviation of
%%     the plotted polygon from the true arc, the 0.651 % at ten points per
%%     decade and the 0.30 % instrument error against which both were read.
%%     The rule is that a number keeps its qualifier or goes with it: the
%%     0.165 % only means anything beside the instrument error, and the
%%     instrument error only means anything beside the density claim it
%%     licenses, so all four go together rather than any one standing bare.
%%     ARCHIVE: ~/chips/castalia/simruns/eisfigs2_20260826, README S1 (and the
%%     0.074 % at 30 points/decade, which was never in the caption).  All 143
%%     points ARE still plotted -- no mark thinning -- see the note above.
%%   - "with the applied potential read at the reference electrode and the
%%     current at the transimpedance output": the sensing architecture, still
%%     stated in the section's own opening paragraph ("A1 applies it at the
%%     reference electrode, A2 and the converter read the resulting current
%%     through the same programmable feedback resistor"), so it is present in
%%     the section, once, and not repeated here.
%%   - the explicit \SIlist of ring frequencies -> "rings mark the decades".
%%     NOT a content cut: every ring carries its own printed label inside the
%%     panels (0.01 Hz / 0.1 Hz / 1 Hz / 10 Hz / 100 Hz nodes), so the list was
%%     the caption reading the figure aloud.
%% WHAT HAD TO STAY, and did: the INDEPENDENT panel scales and the equal axes
%% (without both, a reader misreads the arcs); the agreement figures and where
%% they degrade (0.44 % / 1.03 deg below 20 kHz -> 2.94 % / 4.67 deg at
%% 100 kHz); the gain-code coverage limit (236 ohm - 150.9 kOhm at code 0, so
%% the fresh arm above 151 kOhm is a code-20 read in firmware).
\providecommand{\MaestroRoot}{include/analog/}
\providecolor{mtxA}{HTML}{2A78D6}
\begin{figure}[htbp]
  \centering
  {\pgfplotsset{compat=1.18}%
  \pgfplotsset{
    nyqax/.style={
      scale only axis=true,
      width=3.95cm, height=2.1725cm,
      grid=both,
      major grid style={line width=0.2pt, draw=black!14},
      minor grid style={line width=0.2pt, draw=black!7},
      axis line style={line width=0.4pt, draw=black!45},
      tick style={line width=0.4pt, draw=black!45},
      tick align=outside,
      label style={font=\scriptsize}, tick label style={font=\scriptsize},
      xmin=0, clip=false,
    }}%
  \begin{tikzpicture}
    \begin{axis}[nyqax,
      xmax=1.06106, ymin=-0.0262612, ymax=0.557322,
      xtick={0,0.5,1}, ytick={0,0.2,0.4},
      xlabel={\(Z'\)~[\si{\mega\ohm}]}, ylabel={\(-Z''\)~[\si{\mega\ohm}]},
      title={\footnotesize\textbf{fresh}},
    ]
      \addplot[black!65, line width=0.7pt, no marks]
        table[x expr=\thisrow{x}/1e6, y expr=\thisrow{y}/1e6]
        {\MaestroRoot data/eis_nyq_fresh_true.dat};
      \addplot[only marks, mark=*, mark size=0.85pt, mtxA,
               mark options={fill=mtxA, draw=none}]
        table[x expr=\thisrow{x}/1e6, y expr=\thisrow{y}/1e6]
        {\MaestroRoot data/eis_nyq_fresh_meas.dat};
      \addplot[only marks, mark=o, mark size=1.9pt, black, line width=0.45pt]
        table[x expr=\thisrow{x}/1e6, y expr=\thisrow{y}/1e6]
        {\MaestroRoot data/eis_nyq_fresh_dec.dat};
      \node[left=2pt, font=\tiny] at (axis cs:1.00096,0.00628294) {\SI{0.01}{\hertz}};
      \node[above left=1pt, font=\tiny] at (axis cs:0.997068,0.0625848) {\SI{0.1}{\hertz}};
      \node[above right=1pt, font=\tiny] at (axis cs:0.717957,0.450477) {\SI{1}{\hertz}};
      \node[below right=1pt, font=\tiny] at (axis cs:0.0257045,0.155223) {\SI{10}{\hertz}};
      \node[right=2pt, font=\tiny] at (axis cs:0.00125324,0.0159115) {\SI{100}{\hertz}};
    \end{axis}
  \end{tikzpicture}\hfill
  \begin{tikzpicture}
    \begin{axis}[nyqax,
      xmax=107.06, ymin=-2.64974, ymax=56.2333,
      xtick={0,50,100}, ytick={0,20,40},
      xlabel={\(Z'\)~[\si{\kilo\ohm}]}, ylabel={\(-Z''\)~[\si{\kilo\ohm}]},
      title={\footnotesize\textbf{typical}},
    ]
      \addplot[black!65, line width=0.7pt, no marks]
        table[x expr=\thisrow{x}/1000, y expr=\thisrow{y}/1000]
        {\MaestroRoot data/eis_nyq_typical_true.dat};
      \addplot[only marks, mark=*, mark size=0.85pt, mtxA,
               mark options={fill=mtxA, draw=none}]
        table[x expr=\thisrow{x}/1000, y expr=\thisrow{y}/1000]
        {\MaestroRoot data/eis_nyq_typical_meas.dat};
      \addplot[only marks, mark=o, mark size=1.9pt, black, line width=0.45pt]
        table[x expr=\thisrow{x}/1000, y expr=\thisrow{y}/1000]
        {\MaestroRoot data/eis_nyq_typical_dec.dat};
      \node[left=2pt, font=\tiny] at (axis cs:100.996,0.628294) {\SI{0.01}{\hertz}};
      \node[above left=1pt, font=\tiny] at (axis cs:100.607,6.25848) {\SI{0.1}{\hertz}};
      \node[above right=1pt, font=\tiny] at (axis cs:72.6957,45.0477) {\SI{1}{\hertz}};
      \node[below right=1pt, font=\tiny] at (axis cs:3.47045,15.5223) {\SI{10}{\hertz}};
      \node[right=2pt, font=\tiny] at (axis cs:1.02532,1.59115) {\SI{100}{\hertz}};
    \end{axis}
  \end{tikzpicture}\hfill
  \begin{tikzpicture}
    \begin{axis}[nyqax,
      xmax=11.66, ymin=-0.288585, ymax=6.12441,
      xtick={0,5,10}, ytick={0,2,4,6},
      xlabel={\(Z'\)~[\si{\kilo\ohm}]}, ylabel={\(-Z''\)~[\si{\kilo\ohm}]},
      title={\footnotesize\textbf{fouled}},
    ]
      \addplot[black!65, line width=0.7pt, no marks]
        table[x expr=\thisrow{x}/1000, y expr=\thisrow{y}/1000]
        {\MaestroRoot data/eis_nyq_fouled_true.dat};
      \addplot[only marks, mark=*, mark size=0.85pt, mtxA,
               mark options={fill=mtxA, draw=none}]
        table[x expr=\thisrow{x}/1000, y expr=\thisrow{y}/1000]
        {\MaestroRoot data/eis_nyq_fouled_meas.dat};
      \addplot[only marks, mark=o, mark size=1.9pt, black, line width=0.45pt]
        table[x expr=\thisrow{x}/1000, y expr=\thisrow{y}/1000]
        {\MaestroRoot data/eis_nyq_fouled_dec.dat};
      \node[left=2pt, font=\tiny] at (axis cs:10.9996,0.0628294) {\SI{0.01}{\hertz}};
      \node[above left=1pt, font=\tiny] at (axis cs:10.9607,0.625848) {\SI{0.1}{\hertz}};
      \node[above right=1pt, font=\tiny] at (axis cs:8.16957,4.50477) {\SI{1}{\hertz}};
      \node[below right=1pt, font=\tiny] at (axis cs:1.24705,1.55223) {\SI{10}{\hertz}};
      \node[right=2pt, font=\tiny] at (axis cs:1.00253,0.159115) {\SI{100}{\hertz}};
    \end{axis}
  \end{tikzpicture}}
  \caption[{Nyquist arcs of the three electrode models, measured against
  analytic.}]{Nyquist arcs of the three electrode models measured by the
  channel (dots) against the analytic impedance of the load (line),
  \SI{0.01}{\hertz} to \SI{100}{\kilo\hertz} at \num{20} points per decade,
  gain code 0, nominal process point, \SI{37}{\celsius}; rings mark the
  decades. \textbf{The three panels have independent scales} ---
  \SI{1}{\mega\ohm}, \SI{100}{\kilo\ohm}, \SI{10}{\kilo\ohm} --- and each has
  equal axes, so a departure from a circle is a measurement and not an aspect
  ratio. Agreement is within \SI{0.44}{\percent} of magnitude and
  \SI{1.03}{\degree} of phase below \SI{20}{\kilo\hertz} on all three loads,
  rising to \SI{2.94}{\percent} and \SI{4.67}{\degree} at
  \SI{100}{\kilo\hertz}. Gain code 0 covers \SI{236}{\ohm} to
  \SI{150.9}{\kilo\ohm}, a limit the small-signal sweep does not have, so the
  fresh arm above \SI{151}{\kilo\ohm} is a gain-code-20 read in firmware.}
  \label{fig:eis-nyquist}
\end{figure}

%% back after the conditions paragraph, and put "(Figure~\ref{fig:eis-nyquist})"
%% back on the sentence that now ends "Sweeps carry 20 points per decade."
%% ARCHIVE that holds this figure and every datum in it:
%%   ~/chips/castalia/simruns/eisfigs2_20260826  (build/fig1_nyquist.tex,
%%   data/fig1_{fresh,typical,fouled}_{true,meas,mark}.dat), reduced from
%%   ~/chips/castalia/simruns/eisdeep_20260826 csv/nyquist_*_n0.csv.
%% The caption's facts and the disposition of each are written out fact by fact
%% in note_eis_intro.tex's header, and summarised in note_eis_accuracy.tex's.
%% In short: the agreement figures and the gain-code-0 coverage limit duplicate
%% or are arithmetic on headlines the prose still carries; the panel scales,
%% the equal axes and the "fresh arm above 151 kOhm" note are figure-reading
%% caveats with no figure to read; and the one number that lived only here --
%% 2.94 % and 4.67 deg at 100 kHz -- was CUT rather than moved, because it is
%% the far end of this figure's own code-0 envelope and a decade above the
%% usable band the section states.  It is in the archive above.
%% ======================================================================
%%
%% HAND-WRITTEN fragment -- NOT generated by maestro2tex (raw Spectre runs, not
%% a Maestro results database; there is no maestro2tex config for any *_eis_*
%% fragment).  Precedent for the style block and \MaestroRoot convention:
%% fig_pixtop_noise.tex, fig_eis_error.tex.
%%
%% Drop-in source: ~/chips/castalia/simruns/eisfigs2_20260826/build/
%% fig1_nyquist.tex, with the standalone preamble, the in-figure legend row and
%% the in-figure note block removed -- the note is the caption here.
%% Data: data/eis_nyq_{fresh,typical,fouled}_{true,meas,dec}.dat, copied from
%% eisfigs2_20260826/data/fig1_<load>_{true,meas,mark}.dat, reduced there from
%% ~/chips/castalia/simruns/eisdeep_20260826/csv/nyquist_{fresh,typical,
%% fouled}_n0.csv -- three of that campaign's 30 ac_dense cells, 143
%% frequencies each, 20 points/decade from 0.01 Hz, columns ReZ_meas/ImZ_meas
%% (architecture A) and ReZ_true/ImZ_true (the closed-form Randles impedance)
%% taken as written.  Run provenance: ac_dense/ac_results.json, csv/index.json.
%% Units are ohms in the .dat; the axes divide by 1e6 (fresh) or 1e3.
%%
%% EQUAL AXES ARE LOAD-BEARING.  Every panel has height = 0.55 x width and
%% limits set to the same 0.55 ratio, so the arc is geometrically a semicircle
%% and a departure from a circle is a datum.  Do not change one of width,
%% height or the limits without changing all three.  All 143 measured points
%% are plotted -- the density IS the result, so no mark thinning here (unlike
%% fig_eis_error, where the curve and not the grid is the subject).
%%
%% Key numbers (eisfigs2 README S1): chord deviation of the plotted polygon
%% from the true arc, as a fraction of the arc radius and identical for all
%% three loads because the shape is -- 0.165 % at 20 points/decade, 0.651 % at
%% the archived 10, 0.074 % at 30, against a 0.30 % instrument error.  |D|Z||
%% fresh 0.294 / typical 0.423 / fouled 0.436 % worst case below 20 kHz and
%% 2.808 / 2.925 / 2.936 % at 100 kHz; phase <= 1.031 deg below 20 kHz and
%% <= 4.667 deg at 100 kHz.  All three arcs peak at 1.59 Hz (R_ct C_dl = 0.1 s).
%% Gain code N = 0, R_f = 19.63 kOhm measured in the closed loop.
%% Caveats that had to survive into the caption: the panels have INDEPENDENT
%% scales and equal axes; the ac analysis carries no converter, so the fresh
%% arm above |Z|max = 150.9 kOhm at N = 0 would be taken at N = 20 in firmware
%% (coverage_map.txt).  Nominal corner castalia_typ_25, 37 C, no mismatch -- the
%% corner box on this quantity is 1.6-1.8x (eisdeep S5.3) and the Monte Carlo
%% spread on a stored-law reading is 10 % (~/chips/castalia/simruns/
%% eisdeep_20260826 mc_report/; fig_eis_mccal.tex plotted it until the second
%% cut of 2026-08-26 de-inputted it, and it is still on disk).
%% Not drawn: corner overlays.  The 120 dense corner spectra exist (csv/corner_*)
%% but 5 corners x 3 panels destroys the one thing this figure is for.
%% Library cell names are kept out of the rendered prose.
%% 2026-08-26: the Table~\ref{tab:eis-coverage} pointer was dropped from the
%% caption -- that table went with the cut of note_eis_coverage.tex.  The 236 ohm
%% - 150.9 kOhm window it named is stated inline here and its source is
%% ~/chips/castalia/simruns/eisdemo_20260825/coverage_map.txt.
%% 2026-08-27: the caption was TIGHTENED (wording only -- no number, condition
%% or caveat removed).  Cut: "of them" after all 143; "the loads span" before
%% the three decade scales; "a semicircle reads as a semicircle and" (the
%% second half of the clause says it); "no longer"/"the archived" as filler.
%% "the current measured" became "the current at the transimpedance output",
%% which is the architecture the run used.  Done with the same pass that
%% replaced fig_eis_codelevel with tab_eis_codelevel.
%%
%% 2026-08-27, SECOND PASS -- CONTENT CUT, NOT WORDING.  The earlier pass above
%% was wording only and left the caption at 5 sentences but TWELVE rendered
%% lines (163 pt in the built PDF), which is over the chapter's 2-5 sentence
%% caption style by volume even though it met it by count.  Now 4 sentences.
%% WHAT LEFT THE CAPTION, and where it still lives:
%%   - THE DENSITY CLAIM AND THE CHORD-DEVIATION RESULT, cut together as a
%%     unit: "all 143" of the plotted points, the 0.165 % chord deviation of
%%     the plotted polygon from the true arc, the 0.651 % at ten points per
%%     decade and the 0.30 % instrument error against which both were read.
%%     The rule is that a number keeps its qualifier or goes with it: the
%%     0.165 % only means anything beside the instrument error, and the
%%     instrument error only means anything beside the density claim it
%%     licenses, so all four go together rather than any one standing bare.
%%     ARCHIVE: ~/chips/castalia/simruns/eisfigs2_20260826, README S1 (and the
%%     0.074 % at 30 points/decade, which was never in the caption).  All 143
%%     points ARE still plotted -- no mark thinning -- see the note above.
%%   - "with the applied potential read at the reference electrode and the
%%     current at the transimpedance output": the sensing architecture, still
%%     stated in the section's own opening paragraph ("A1 applies it at the
%%     reference electrode, A2 and the converter read the resulting current
%%     through the same programmable feedback resistor"), so it is present in
%%     the section, once, and not repeated here.
%%   - the explicit \SIlist of ring frequencies -> "rings mark the decades".
%%     NOT a content cut: every ring carries its own printed label inside the
%%     panels (0.01 Hz / 0.1 Hz / 1 Hz / 10 Hz / 100 Hz nodes), so the list was
%%     the caption reading the figure aloud.
%% WHAT HAD TO STAY, and did: the INDEPENDENT panel scales and the equal axes
%% (without both, a reader misreads the arcs); the agreement figures and where
%% they degrade (0.44 % / 1.03 deg below 20 kHz -> 2.94 % / 4.67 deg at
%% 100 kHz); the gain-code coverage limit (236 ohm - 150.9 kOhm at code 0, so
%% the fresh arm above 151 kOhm is a code-20 read in firmware).
\providecommand{\MaestroRoot}{include/analog/}
\providecolor{mtxA}{HTML}{2A78D6}
\begin{figure}[htbp]
  \centering
  {\pgfplotsset{compat=1.18}%
  \pgfplotsset{
    nyqax/.style={
      scale only axis=true,
      width=3.95cm, height=2.1725cm,
      grid=both,
      major grid style={line width=0.2pt, draw=black!14},
      minor grid style={line width=0.2pt, draw=black!7},
      axis line style={line width=0.4pt, draw=black!45},
      tick style={line width=0.4pt, draw=black!45},
      tick align=outside,
      label style={font=\scriptsize}, tick label style={font=\scriptsize},
      xmin=0, clip=false,
    }}%
  \begin{tikzpicture}
    \begin{axis}[nyqax,
      xmax=1.06106, ymin=-0.0262612, ymax=0.557322,
      xtick={0,0.5,1}, ytick={0,0.2,0.4},
      xlabel={\(Z'\)~[\si{\mega\ohm}]}, ylabel={\(-Z''\)~[\si{\mega\ohm}]},
      title={\footnotesize\textbf{fresh}},
    ]
      \addplot[black!65, line width=0.7pt, no marks]
        table[x expr=\thisrow{x}/1e6, y expr=\thisrow{y}/1e6]
        {\MaestroRoot data/eis_nyq_fresh_true.dat};
      \addplot[only marks, mark=*, mark size=0.85pt, mtxA,
               mark options={fill=mtxA, draw=none}]
        table[x expr=\thisrow{x}/1e6, y expr=\thisrow{y}/1e6]
        {\MaestroRoot data/eis_nyq_fresh_meas.dat};
      \addplot[only marks, mark=o, mark size=1.9pt, black, line width=0.45pt]
        table[x expr=\thisrow{x}/1e6, y expr=\thisrow{y}/1e6]
        {\MaestroRoot data/eis_nyq_fresh_dec.dat};
      \node[left=2pt, font=\tiny] at (axis cs:1.00096,0.00628294) {\SI{0.01}{\hertz}};
      \node[above left=1pt, font=\tiny] at (axis cs:0.997068,0.0625848) {\SI{0.1}{\hertz}};
      \node[above right=1pt, font=\tiny] at (axis cs:0.717957,0.450477) {\SI{1}{\hertz}};
      \node[below right=1pt, font=\tiny] at (axis cs:0.0257045,0.155223) {\SI{10}{\hertz}};
      \node[right=2pt, font=\tiny] at (axis cs:0.00125324,0.0159115) {\SI{100}{\hertz}};
    \end{axis}
  \end{tikzpicture}\hfill
  \begin{tikzpicture}
    \begin{axis}[nyqax,
      xmax=107.06, ymin=-2.64974, ymax=56.2333,
      xtick={0,50,100}, ytick={0,20,40},
      xlabel={\(Z'\)~[\si{\kilo\ohm}]}, ylabel={\(-Z''\)~[\si{\kilo\ohm}]},
      title={\footnotesize\textbf{typical}},
    ]
      \addplot[black!65, line width=0.7pt, no marks]
        table[x expr=\thisrow{x}/1000, y expr=\thisrow{y}/1000]
        {\MaestroRoot data/eis_nyq_typical_true.dat};
      \addplot[only marks, mark=*, mark size=0.85pt, mtxA,
               mark options={fill=mtxA, draw=none}]
        table[x expr=\thisrow{x}/1000, y expr=\thisrow{y}/1000]
        {\MaestroRoot data/eis_nyq_typical_meas.dat};
      \addplot[only marks, mark=o, mark size=1.9pt, black, line width=0.45pt]
        table[x expr=\thisrow{x}/1000, y expr=\thisrow{y}/1000]
        {\MaestroRoot data/eis_nyq_typical_dec.dat};
      \node[left=2pt, font=\tiny] at (axis cs:100.996,0.628294) {\SI{0.01}{\hertz}};
      \node[above left=1pt, font=\tiny] at (axis cs:100.607,6.25848) {\SI{0.1}{\hertz}};
      \node[above right=1pt, font=\tiny] at (axis cs:72.6957,45.0477) {\SI{1}{\hertz}};
      \node[below right=1pt, font=\tiny] at (axis cs:3.47045,15.5223) {\SI{10}{\hertz}};
      \node[right=2pt, font=\tiny] at (axis cs:1.02532,1.59115) {\SI{100}{\hertz}};
    \end{axis}
  \end{tikzpicture}\hfill
  \begin{tikzpicture}
    \begin{axis}[nyqax,
      xmax=11.66, ymin=-0.288585, ymax=6.12441,
      xtick={0,5,10}, ytick={0,2,4,6},
      xlabel={\(Z'\)~[\si{\kilo\ohm}]}, ylabel={\(-Z''\)~[\si{\kilo\ohm}]},
      title={\footnotesize\textbf{fouled}},
    ]
      \addplot[black!65, line width=0.7pt, no marks]
        table[x expr=\thisrow{x}/1000, y expr=\thisrow{y}/1000]
        {\MaestroRoot data/eis_nyq_fouled_true.dat};
      \addplot[only marks, mark=*, mark size=0.85pt, mtxA,
               mark options={fill=mtxA, draw=none}]
        table[x expr=\thisrow{x}/1000, y expr=\thisrow{y}/1000]
        {\MaestroRoot data/eis_nyq_fouled_meas.dat};
      \addplot[only marks, mark=o, mark size=1.9pt, black, line width=0.45pt]
        table[x expr=\thisrow{x}/1000, y expr=\thisrow{y}/1000]
        {\MaestroRoot data/eis_nyq_fouled_dec.dat};
      \node[left=2pt, font=\tiny] at (axis cs:10.9996,0.0628294) {\SI{0.01}{\hertz}};
      \node[above left=1pt, font=\tiny] at (axis cs:10.9607,0.625848) {\SI{0.1}{\hertz}};
      \node[above right=1pt, font=\tiny] at (axis cs:8.16957,4.50477) {\SI{1}{\hertz}};
      \node[below right=1pt, font=\tiny] at (axis cs:1.24705,1.55223) {\SI{10}{\hertz}};
      \node[right=2pt, font=\tiny] at (axis cs:1.00253,0.159115) {\SI{100}{\hertz}};
    \end{axis}
  \end{tikzpicture}}
  \caption[{Nyquist arcs of the three electrode models, measured against
  analytic.}]{Nyquist arcs of the three electrode models measured by the
  channel (dots) against the analytic impedance of the load (line),
  \SI{0.01}{\hertz} to \SI{100}{\kilo\hertz} at \num{20} points per decade,
  gain code 0, nominal process point, \SI{37}{\celsius}; rings mark the
  decades. \textbf{The three panels have independent scales} ---
  \SI{1}{\mega\ohm}, \SI{100}{\kilo\ohm}, \SI{10}{\kilo\ohm} --- and each has
  equal axes, so a departure from a circle is a measurement and not an aspect
  ratio. Agreement is within \SI{0.44}{\percent} of magnitude and
  \SI{1.03}{\degree} of phase below \SI{20}{\kilo\hertz} on all three loads,
  rising to \SI{2.94}{\percent} and \SI{4.67}{\degree} at
  \SI{100}{\kilo\hertz}. Gain code 0 covers \SI{236}{\ohm} to
  \SI{150.9}{\kilo\ohm}, a limit the small-signal sweep does not have, so the
  fresh arm above \SI{151}{\kilo\ohm} is a gain-code-20 read in firmware.}
  \label{fig:eis-nyquist}
\end{figure}

%% back after the conditions paragraph, and put "(Figure~\ref{fig:eis-nyquist})"
%% back on the sentence that now ends "Sweeps carry 20 points per decade."
%% ARCHIVE that holds this figure and every datum in it:
%%   ~/chips/castalia/simruns/eisfigs2_20260826  (build/fig1_nyquist.tex,
%%   data/fig1_{fresh,typical,fouled}_{true,meas,mark}.dat), reduced from
%%   ~/chips/castalia/simruns/eisdeep_20260826 csv/nyquist_*_n0.csv.
%% The caption's facts and the disposition of each are written out fact by fact
%% in note_eis_intro.tex's header, and summarised in note_eis_accuracy.tex's.
%% In short: the agreement figures and the gain-code-0 coverage limit duplicate
%% or are arithmetic on headlines the prose still carries; the panel scales,
%% the equal axes and the "fresh arm above 151 kOhm" note are figure-reading
%% caveats with no figure to read; and the one number that lived only here --
%% 2.94 % and 4.67 deg at 100 kHz -- was CUT rather than moved, because it is
%% the far end of this figure's own code-0 envelope and a decade above the
%% usable band the section states.  It is in the archive above.
%% ======================================================================
%%
%% HAND-WRITTEN fragment -- NOT generated by maestro2tex (raw Spectre runs, not
%% a Maestro results database; there is no maestro2tex config for any *_eis_*
%% fragment).  Precedent for the style block and \MaestroRoot convention:
%% fig_pixtop_noise.tex, fig_eis_error.tex.
%%
%% Drop-in source: ~/chips/castalia/simruns/eisfigs2_20260826/build/
%% fig1_nyquist.tex, with the standalone preamble, the in-figure legend row and
%% the in-figure note block removed -- the note is the caption here.
%% Data: data/eis_nyq_{fresh,typical,fouled}_{true,meas,dec}.dat, copied from
%% eisfigs2_20260826/data/fig1_<load>_{true,meas,mark}.dat, reduced there from
%% ~/chips/castalia/simruns/eisdeep_20260826/csv/nyquist_{fresh,typical,
%% fouled}_n0.csv -- three of that campaign's 30 ac_dense cells, 143
%% frequencies each, 20 points/decade from 0.01 Hz, columns ReZ_meas/ImZ_meas
%% (architecture A) and ReZ_true/ImZ_true (the closed-form Randles impedance)
%% taken as written.  Run provenance: ac_dense/ac_results.json, csv/index.json.
%% Units are ohms in the .dat; the axes divide by 1e6 (fresh) or 1e3.
%%
%% EQUAL AXES ARE LOAD-BEARING.  Every panel has height = 0.55 x width and
%% limits set to the same 0.55 ratio, so the arc is geometrically a semicircle
%% and a departure from a circle is a datum.  Do not change one of width,
%% height or the limits without changing all three.  All 143 measured points
%% are plotted -- the density IS the result, so no mark thinning here (unlike
%% fig_eis_error, where the curve and not the grid is the subject).
%%
%% Key numbers (eisfigs2 README S1): chord deviation of the plotted polygon
%% from the true arc, as a fraction of the arc radius and identical for all
%% three loads because the shape is -- 0.165 % at 20 points/decade, 0.651 % at
%% the archived 10, 0.074 % at 30, against a 0.30 % instrument error.  |D|Z||
%% fresh 0.294 / typical 0.423 / fouled 0.436 % worst case below 20 kHz and
%% 2.808 / 2.925 / 2.936 % at 100 kHz; phase <= 1.031 deg below 20 kHz and
%% <= 4.667 deg at 100 kHz.  All three arcs peak at 1.59 Hz (R_ct C_dl = 0.1 s).
%% Gain code N = 0, R_f = 19.63 kOhm measured in the closed loop.
%% Caveats that had to survive into the caption: the panels have INDEPENDENT
%% scales and equal axes; the ac analysis carries no converter, so the fresh
%% arm above |Z|max = 150.9 kOhm at N = 0 would be taken at N = 20 in firmware
%% (coverage_map.txt).  Nominal corner castalia_typ_25, 37 C, no mismatch -- the
%% corner box on this quantity is 1.6-1.8x (eisdeep S5.3) and the Monte Carlo
%% spread on a stored-law reading is 10 % (~/chips/castalia/simruns/
%% eisdeep_20260826 mc_report/; fig_eis_mccal.tex plotted it until the second
%% cut of 2026-08-26 de-inputted it, and it is still on disk).
%% Not drawn: corner overlays.  The 120 dense corner spectra exist (csv/corner_*)
%% but 5 corners x 3 panels destroys the one thing this figure is for.
%% Library cell names are kept out of the rendered prose.
%% 2026-08-26: the Table~\ref{tab:eis-coverage} pointer was dropped from the
%% caption -- that table went with the cut of note_eis_coverage.tex.  The 236 ohm
%% - 150.9 kOhm window it named is stated inline here and its source is
%% ~/chips/castalia/simruns/eisdemo_20260825/coverage_map.txt.
%% 2026-08-27: the caption was TIGHTENED (wording only -- no number, condition
%% or caveat removed).  Cut: "of them" after all 143; "the loads span" before
%% the three decade scales; "a semicircle reads as a semicircle and" (the
%% second half of the clause says it); "no longer"/"the archived" as filler.
%% "the current measured" became "the current at the transimpedance output",
%% which is the architecture the run used.  Done with the same pass that
%% replaced fig_eis_codelevel with tab_eis_codelevel.
%%
%% 2026-08-27, SECOND PASS -- CONTENT CUT, NOT WORDING.  The earlier pass above
%% was wording only and left the caption at 5 sentences but TWELVE rendered
%% lines (163 pt in the built PDF), which is over the chapter's 2-5 sentence
%% caption style by volume even though it met it by count.  Now 4 sentences.
%% WHAT LEFT THE CAPTION, and where it still lives:
%%   - THE DENSITY CLAIM AND THE CHORD-DEVIATION RESULT, cut together as a
%%     unit: "all 143" of the plotted points, the 0.165 % chord deviation of
%%     the plotted polygon from the true arc, the 0.651 % at ten points per
%%     decade and the 0.30 % instrument error against which both were read.
%%     The rule is that a number keeps its qualifier or goes with it: the
%%     0.165 % only means anything beside the instrument error, and the
%%     instrument error only means anything beside the density claim it
%%     licenses, so all four go together rather than any one standing bare.
%%     ARCHIVE: ~/chips/castalia/simruns/eisfigs2_20260826, README S1 (and the
%%     0.074 % at 30 points/decade, which was never in the caption).  All 143
%%     points ARE still plotted -- no mark thinning -- see the note above.
%%   - "with the applied potential read at the reference electrode and the
%%     current at the transimpedance output": the sensing architecture, still
%%     stated in the section's own opening paragraph ("A1 applies it at the
%%     reference electrode, A2 and the converter read the resulting current
%%     through the same programmable feedback resistor"), so it is present in
%%     the section, once, and not repeated here.
%%   - the explicit \SIlist of ring frequencies -> "rings mark the decades".
%%     NOT a content cut: every ring carries its own printed label inside the
%%     panels (0.01 Hz / 0.1 Hz / 1 Hz / 10 Hz / 100 Hz nodes), so the list was
%%     the caption reading the figure aloud.
%% WHAT HAD TO STAY, and did: the INDEPENDENT panel scales and the equal axes
%% (without both, a reader misreads the arcs); the agreement figures and where
%% they degrade (0.44 % / 1.03 deg below 20 kHz -> 2.94 % / 4.67 deg at
%% 100 kHz); the gain-code coverage limit (236 ohm - 150.9 kOhm at code 0, so
%% the fresh arm above 151 kOhm is a code-20 read in firmware).
\providecommand{\MaestroRoot}{include/analog/}
\providecolor{mtxA}{HTML}{2A78D6}
\begin{figure}[htbp]
  \centering
  {\pgfplotsset{compat=1.18}%
  \pgfplotsset{
    nyqax/.style={
      scale only axis=true,
      width=3.95cm, height=2.1725cm,
      grid=both,
      major grid style={line width=0.2pt, draw=black!14},
      minor grid style={line width=0.2pt, draw=black!7},
      axis line style={line width=0.4pt, draw=black!45},
      tick style={line width=0.4pt, draw=black!45},
      tick align=outside,
      label style={font=\scriptsize}, tick label style={font=\scriptsize},
      xmin=0, clip=false,
    }}%
  \begin{tikzpicture}
    \begin{axis}[nyqax,
      xmax=1.06106, ymin=-0.0262612, ymax=0.557322,
      xtick={0,0.5,1}, ytick={0,0.2,0.4},
      xlabel={\(Z'\)~[\si{\mega\ohm}]}, ylabel={\(-Z''\)~[\si{\mega\ohm}]},
      title={\footnotesize\textbf{fresh}},
    ]
      \addplot[black!65, line width=0.7pt, no marks]
        table[x expr=\thisrow{x}/1e6, y expr=\thisrow{y}/1e6]
        {\MaestroRoot data/eis_nyq_fresh_true.dat};
      \addplot[only marks, mark=*, mark size=0.85pt, mtxA,
               mark options={fill=mtxA, draw=none}]
        table[x expr=\thisrow{x}/1e6, y expr=\thisrow{y}/1e6]
        {\MaestroRoot data/eis_nyq_fresh_meas.dat};
      \addplot[only marks, mark=o, mark size=1.9pt, black, line width=0.45pt]
        table[x expr=\thisrow{x}/1e6, y expr=\thisrow{y}/1e6]
        {\MaestroRoot data/eis_nyq_fresh_dec.dat};
      \node[left=2pt, font=\tiny] at (axis cs:1.00096,0.00628294) {\SI{0.01}{\hertz}};
      \node[above left=1pt, font=\tiny] at (axis cs:0.997068,0.0625848) {\SI{0.1}{\hertz}};
      \node[above right=1pt, font=\tiny] at (axis cs:0.717957,0.450477) {\SI{1}{\hertz}};
      \node[below right=1pt, font=\tiny] at (axis cs:0.0257045,0.155223) {\SI{10}{\hertz}};
      \node[right=2pt, font=\tiny] at (axis cs:0.00125324,0.0159115) {\SI{100}{\hertz}};
    \end{axis}
  \end{tikzpicture}\hfill
  \begin{tikzpicture}
    \begin{axis}[nyqax,
      xmax=107.06, ymin=-2.64974, ymax=56.2333,
      xtick={0,50,100}, ytick={0,20,40},
      xlabel={\(Z'\)~[\si{\kilo\ohm}]}, ylabel={\(-Z''\)~[\si{\kilo\ohm}]},
      title={\footnotesize\textbf{typical}},
    ]
      \addplot[black!65, line width=0.7pt, no marks]
        table[x expr=\thisrow{x}/1000, y expr=\thisrow{y}/1000]
        {\MaestroRoot data/eis_nyq_typical_true.dat};
      \addplot[only marks, mark=*, mark size=0.85pt, mtxA,
               mark options={fill=mtxA, draw=none}]
        table[x expr=\thisrow{x}/1000, y expr=\thisrow{y}/1000]
        {\MaestroRoot data/eis_nyq_typical_meas.dat};
      \addplot[only marks, mark=o, mark size=1.9pt, black, line width=0.45pt]
        table[x expr=\thisrow{x}/1000, y expr=\thisrow{y}/1000]
        {\MaestroRoot data/eis_nyq_typical_dec.dat};
      \node[left=2pt, font=\tiny] at (axis cs:100.996,0.628294) {\SI{0.01}{\hertz}};
      \node[above left=1pt, font=\tiny] at (axis cs:100.607,6.25848) {\SI{0.1}{\hertz}};
      \node[above right=1pt, font=\tiny] at (axis cs:72.6957,45.0477) {\SI{1}{\hertz}};
      \node[below right=1pt, font=\tiny] at (axis cs:3.47045,15.5223) {\SI{10}{\hertz}};
      \node[right=2pt, font=\tiny] at (axis cs:1.02532,1.59115) {\SI{100}{\hertz}};
    \end{axis}
  \end{tikzpicture}\hfill
  \begin{tikzpicture}
    \begin{axis}[nyqax,
      xmax=11.66, ymin=-0.288585, ymax=6.12441,
      xtick={0,5,10}, ytick={0,2,4,6},
      xlabel={\(Z'\)~[\si{\kilo\ohm}]}, ylabel={\(-Z''\)~[\si{\kilo\ohm}]},
      title={\footnotesize\textbf{fouled}},
    ]
      \addplot[black!65, line width=0.7pt, no marks]
        table[x expr=\thisrow{x}/1000, y expr=\thisrow{y}/1000]
        {\MaestroRoot data/eis_nyq_fouled_true.dat};
      \addplot[only marks, mark=*, mark size=0.85pt, mtxA,
               mark options={fill=mtxA, draw=none}]
        table[x expr=\thisrow{x}/1000, y expr=\thisrow{y}/1000]
        {\MaestroRoot data/eis_nyq_fouled_meas.dat};
      \addplot[only marks, mark=o, mark size=1.9pt, black, line width=0.45pt]
        table[x expr=\thisrow{x}/1000, y expr=\thisrow{y}/1000]
        {\MaestroRoot data/eis_nyq_fouled_dec.dat};
      \node[left=2pt, font=\tiny] at (axis cs:10.9996,0.0628294) {\SI{0.01}{\hertz}};
      \node[above left=1pt, font=\tiny] at (axis cs:10.9607,0.625848) {\SI{0.1}{\hertz}};
      \node[above right=1pt, font=\tiny] at (axis cs:8.16957,4.50477) {\SI{1}{\hertz}};
      \node[below right=1pt, font=\tiny] at (axis cs:1.24705,1.55223) {\SI{10}{\hertz}};
      \node[right=2pt, font=\tiny] at (axis cs:1.00253,0.159115) {\SI{100}{\hertz}};
    \end{axis}
  \end{tikzpicture}}
  \caption[{Nyquist arcs of the three electrode models, measured against
  analytic.}]{Nyquist arcs of the three electrode models measured by the
  channel (dots) against the analytic impedance of the load (line),
  \SI{0.01}{\hertz} to \SI{100}{\kilo\hertz} at \num{20} points per decade,
  gain code 0, nominal process point, \SI{37}{\celsius}; rings mark the
  decades. \textbf{The three panels have independent scales} ---
  \SI{1}{\mega\ohm}, \SI{100}{\kilo\ohm}, \SI{10}{\kilo\ohm} --- and each has
  equal axes, so a departure from a circle is a measurement and not an aspect
  ratio. Agreement is within \SI{0.44}{\percent} of magnitude and
  \SI{1.03}{\degree} of phase below \SI{20}{\kilo\hertz} on all three loads,
  rising to \SI{2.94}{\percent} and \SI{4.67}{\degree} at
  \SI{100}{\kilo\hertz}. Gain code 0 covers \SI{236}{\ohm} to
  \SI{150.9}{\kilo\ohm}, a limit the small-signal sweep does not have, so the
  fresh arm above \SI{151}{\kilo\ohm} is a gain-code-20 read in firmware.}
  \label{fig:eis-nyquist}
\end{figure}
